% Generated draft: requires mathematical and source review.
\documentclass[11pt]{article}
\usepackage[margin=25mm]{geometry}
\usepackage{fontspec}
\setmainfont{Latin Modern Roman}
\usepackage{amsmath,amssymb,mathtools,unicode-math}
\setmathfont{Latin Modern Math}
\usepackage{enumitem,needspace,xurl,hyperref}
\setlength{\parindent}{0pt}
\setlength{\parskip}{5pt}
\setlength{\emergencystretch}{3em}
\newcommand{\N}{\mathbb{N}}
\newcommand{\Z}{\mathbb{Z}}
\newcommand{\Q}{\mathbb{Q}}
\newcommand{\R}{\mathbb{R}}
\newcommand{\C}{\mathbb{C}}
\newcommand{\F}{\mathbb{F}}
\newcommand{\A}{\mathbb{A}}
\newcommand{\PP}{\mathbb{P}}
\newcommand{\E}{\mathbb{E}}
\newcommand{\eps}{\varepsilon}
\newcommand{\dd}{\,\mathrm{d}}
\begin{document}
\section*{UnsolvedMath ID: 701422}
\textbf{Source problem code:} catalogue entry
\subsection{Definitions and mathematical statement}
Put \(\kappa=\aleph_\omega=\sup_{n<\omega}\aleph_n\). Let \(G=(\R,\cdot,e)\) be a group whose underlying set is exactly \(\R\); no commutativity assumption is supplied. A group is free if it has a basis such that every element has a unique reduced word in the basis elements and their inverses. Using the usual convention, call \(G\) \(\kappa\)-free if every subgroup \(H\leq G\) with \(|H|<\kappa\) is a free group.

The source asks, separately for each of the following meanings of ``nicely definable,'' whether
\[
G\text{ is }\aleph_\omega\text{-free}
\quad\Longrightarrow\quad
G\text{ is free}:
\]
\begin{enumerate}
\item[(a)] \(F_\sigma\);
\item[(b)] Borel;
\item[(c)] analytic;
\item[(d)] projective, followed in the supplied source by malformed notation involving \(\mathbb L\) and \(\mathbb R\).
\end{enumerate}
The source does not specify which relation or coding of the group must have the indicated descriptive complexity. Since the underlying set is already \(\R\), its complexity alone does not specify a definability restriction on the group operation.

One possible interpretation of parts (a)--(c) assigns the complexity to the multiplication graph
\[
M_G=\{(x,y,z)\in\R^3:x\cdot y=z\},
\]
with \(\R^3\) carrying its usual product topology. Here an \(F_\sigma\) set is a countable union of closed sets, a Borel set belongs to the sigma-algebra generated by open sets, and an analytic set is a projection of a Borel subset of \(\R^3\times\N^{\N}\), where \(\N^{\N}\) has its usual product topology. Under this explicitly conditional interpretation, each of (a)--(c) asks whether, for every group structure \(G\) on \(\R\),
\[
\left(M_G\in\Gamma\ \text{ and }\
\forall H\leq G\;\bigl(|H|<\aleph_\omega\Longrightarrow H\text{ is free}\bigr)\right)
\Longrightarrow G\text{ is free},
\]
where \(\Gamma\) is respectively the class of \(F_\sigma\), Borel, or analytic subsets of \(\R^3\). The supplied record does not establish that this graph interpretation is intended. Part (d) cannot be given a precise definability predicate from the supplied malformed text; its intended notation, ambient universe, and relationship to projectivity require clarification. The source refers to the paper linked in the Sources for justification of its restrictions, but the relevant qualifications are not reproduced in the record.
\subsection{Short English statement}
Shelah asks whether every \(\aleph_\omega\)-free group with underlying set \(\R\) must be free under each of four proposed definability restrictions. The intended coding of definability is unspecified, and the fourth restriction has malformed notation. The multiplication-graph formulation above is therefore conditional.
\subsection{Sources}
\begin{itemize}
\item Supplied local UnsolvedMath registry record: permanent ID \(701422\), exact source problem code \texttt{catalogue entry}, title ``Kourovka Notebook Problem 17.121.'' The record attributes the problem to S. Shelah, proposed in 2010, and labels it open; this label is supplied registry metadata, not an independently verified assessment.
\item Supplied authoritative-source attribution: \emph{The Kourovka Notebook (Unsolved Problems in Group Theory)}, 17th issue (2010), Problem 17.121; also identified in the record with the 21st issue edition (2026), arXiv:1401.0300v48. The collection description names E. I. Khukhro and V. D. Mazurov as editors. These bibliographic details are supplied attribution and have not been independently checked.
\item The statement supplies \url{https://arxiv.org/pdf/math/0212250.pdf} for ``justification of the restrictions.'' No title, authors, or relevant passage from that paper is supplied, and its contents have not been consulted here.
\end{itemize}
\end{document}
